\subsection{History of Scholarly Search}
\label{sec:history}

Scholarly retrieval has navigated a recurring tension between structured
indexing (controlled vocabularies and typed metadata) and free-text retrieval
(full-content indexing), with each era layering onto the last rather than
swinging between them. Pre-digital catalogues and the early digital
bibliographic databases (MEDLARS, INSPEC, PsycINFO) were anchored to
controlled vocabularies and Boolean intermediaries, whose dominant failure
mode Lancaster diagnosed as vocabulary mismatch rather than algorithmic
weakness~\cite{lancaster1969medlars}. The web era sidestepped that mismatch by
indexing full text at scale through Google Scholar, Scopus, and BM25- and
PageRank-driven
ranking~\cite{robertson1994bm25,brin1998pagerank,falagas2008comparison}, at
the cost of the disciplined precision that controlled vocabularies had
guaranteed. Modern infrastructure (Semantic Scholar, OpenAlex,
OpenCitations) partially reintroduced typed metadata as a navigable layer over
web-scale
text~\cite{ammar2018construction,priem2022openalex,peroni2020opencitations},
and the post-2022 generation of LLM- and embedding-driven tools now sits on
top of those graphs. Each generation has therefore inherited the substrate
of the last, addressed its dominant limitation, and introduced a new one. The
methods examined in Section~\ref{sec:recent-advances}, namely dense
retrieval, graph augmentation, natural-language-to-structured-query
translation, and query reformulation, are best read as the field's current
attempt to reconcile these two paradigms rather than choose between them,
and this thesis takes that reconciliation as its starting point.